\chapter{2019年北京卷（理）}

\section{单选题}

\begin{problem}
已知复数 \(z=2+\i\)，则 \(z\overline z=\)（\quad）.

\choices
    {\(\sqrt3\)}
    {\(\sqrt5\)}
    {\(3\)}
    {\(5\)}
\end{problem}



\begin{problem}
执行如图所示的程序框图，输出的 \(s\) 值为

\begin{center}
\texfigure[max width=6.0cm]{img_repaint/2019/beijing_science/q2_fig1.tex}
\end{center}

\choices
    {\(1\)}
    {\(2\)}
    {\(3\)}
    {\(4\)}
\end{problem}



\begin{problem}
已知直线 \(l\) 的参数方程为
\(\begin{cases}
x=1+3t,\\
y=2+4t,
\end{cases}\)
（其中 \(t\) 为参数），则点 \((1,0)\) 到直线 \(l\) 的距离是（\quad）.

\choices
    {\(\dfrac15\)}
    {\(\dfrac25\)}
    {\(\dfrac45\)}
    {\(\dfrac65\)}
\end{problem}



\begin{problem}
已知椭圆 \(\dfrac{x^2}{a^2}+\dfrac{y^2}{b^2}=1\)（\(a>b>0\)）的离心率为 \(\dfrac12\)，则（\quad）.

\choices
    {\(a^2=2b^2\)}
    {\(3a^2=4b^2\)}
    {\(a=2b\)}
    {\(3a=4b\)}
\end{problem}



\begin{problem}
若 \(x\)，\(y\) 满足 \(\left|x\right|\le1-y\)，且 \(y\ge-1\)，则 \(3x+y\) 的最大值为（\quad）.

\choices
    {\(-7\)}
    {\(1\)}
    {\(5\)}
    {\(7\)}
\end{problem}



\begin{problem}
在天文学中，天体的明暗程度可以用星等或亮度来描述. 两颗星的星等与亮度满足 \(m_2-m_1=\dfrac52\lg\dfrac{E_1}{E_2}\)，其中星等为 \(m_k\) 的星的亮度为 \(E_k\)（\(k=1\)，\(2\)）. 已知太阳的星等是 \(-26.7\)，天狼星的星等是 \(-1.45\)，则太阳与天狼星的亮度的比值为（\quad）.

\choices
    {\(10^{10.1}\)}
    {\(10.1\)}
    {\(\lg10.1\)}
    {\(10^{-10.1}\)}
\end{problem}



\begin{problem}
设点 \(A\)，\(B\)，\(C\) 不共线，则“\(\overrightarrow{AB}\) 与 \(\overrightarrow{AC}\) 的夹角是锐角”是“\(\left|\overrightarrow{AB}+\overrightarrow{AC}\right|>\left|\overrightarrow{BC}\right|\)”的（\quad）.

\choices
    {充分而不必要条件}
    {必要而不充分条件}
    {充分必要条件}
    {既不充分也不必要条件}
\end{problem}



\begin{problem}
数学中有许多形状优美，寓意美好的曲线，曲线 \(C:x^2+y^2=1+\left|x\right|y\) 就是其中之一（如图）. 给出下列三个结论：

\bitmapfigure[max width=6.0cm]{img/2019/beijing_science/q8_fig1.png}

\begin{circlelist}
\item 曲线 \(C\) 恰好经过 \(6\) 个整点（即横，纵坐标均为整数的点）；

\item 曲线 \(C\) 上任意一点到原点的距离都不超过 \(\sqrt2\)；

\item 曲线 \(C\) 所围成的“心形”区域的面积小于 \(3\).
\end{circlelist}
其中，所有正确结论的序号是（\quad）.

\choices
    {\circled{1}}
    {\circled{2}}
    {\circled{1}\circled{2}}
    {\circled{1}\circled{2}\circled{3}}
\end{problem}


\section{填空题}

\begin{problem}
函数 \(\sin^2 2x\) 的最小正周期是\fillinblank{}.
\end{problem}



\begin{problem}
设等差数列 \(\{a_n\}\) 的前 \(n\) 项和为 \(S_n\). 若 \(a_2=-3\)，\(S_5=-10\)，则 \(a_3=\fillinblank{}\)，\(S_n\) 的最小值为\fillinblank{}.
\end{problem}



\begin{problem}
某几何体是由一个正方体去掉一个四棱柱所得，其三视图如图所示. 如果网格纸上小正方形的边长为 \(1\)，那么该几何体的体积为\fillinblank{}.

\bitmapfigure[max width=0.58\linewidth]{img/2019/beijing_science/q11_fig1.png}
\end{problem}



\begin{problem}
已知 \(l\)，\(m\) 是平面 \(\alpha\) 外的两条不同直线. 给出下列三个论断：

\begin{circlelist}
\item \(l\perp m\)；

\item \(m//\alpha\)；

\item \(l\perp\alpha\).
\end{circlelist}
以其中的两个论断作为条件，余下的一个论断作为结论，写出一个正确的命题：\fillinblank{}.
\end{problem}



\begin{problem}
设函数 \(f(x)=\e^x+a\e^{-x}\)（\(a\) 为常数），若 \(f(x)\) 为奇函数，则 \(a=\fillinblank{}\)；若 \(f(x)\) 是 \(\R\) 上的增函数，则 \(a\) 的取值范围是\fillinblank{}.
\end{problem}



\begin{problem}
李明自主创业，在网上经营一家水果店，销售的水果中有草莓，京白梨，西瓜，桃. 价格依次为 \(60\) 元/盒，\(65\) 元/盒，\(80\) 元/盒，\(90\) 元/盒. 为增加销量，李明对这四种水果进行促销：一次购买水果的总价达到 \(120\) 元，顾客就少付 \(x\) 元. 每笔订单顾客网上支付成功后，李明会得到支付款的 \(80\%\).

\begin{enumerate}
\item 当 \(x=10\) 时，顾客一次购买草莓和西瓜各 \(1\) 盒，需要支付\fillinblank{} 元；

\item 在促销活动中，为保证李明每笔订单得到的金额均不低于促销前总价的七折，则 \(x\) 的最大值为\fillinblank{}.
\end{enumerate}
\end{problem}


\section{解答题}

\begin{problem}
在 \(\triangle ABC\) 中，\(a=3\)，\(b-c=2\)，\(\cos B=-\dfrac12\).

\begin{enumerate}
\item 求 \(b\)，\(c\) 的值；

\item 求 \(\sin(B-C)\) 的值.
\end{enumerate}
\end{problem}

\begin{problem}
如图，在四棱锥 \(P\)-\(ABCD\) 中，\(PA\perp\) 平面 \(ABCD\)，\(AD\perp CD\)，\(AD\parallel BC\)，\(PA=AD=CD=2\)，\(BC=3\). \(E\) 为 \(PD\) 的中点，点 \(F\) 在 \(PC\) 上，且 \(\dfrac{PF}{PC}=\dfrac13\).

\begin{enumerate}
\item 求证：\(CD\perp\) 平面 \(PAD\)；

\item 求二面角 \(F\)-\(AE\)-\(P\) 的余弦值；

\item 设点 \(G\) 在 \(PB\) 上，且 \(\dfrac{PG}{PB}=\dfrac23\). 判断直线 \(AG\) 是否在平面 \(AEF\) 内，说明理由.
\end{enumerate}

\bitmapfigure[max width=6.0cm]{img/2019/beijing_science/q16_fig1.png}
\end{problem}


\begin{problem}
改革开放以来，人们的支付方式发生了巨大转变. 近年来，移动支付已成为主要支付方式之一. 为了解某校学生上个月 \(A\)，\(B\) 两种移动支付方式的使用情况，从全校学生中随机抽取了 \(100\) 人，发现样本中 \(A\)，\(B\) 两种支付方式都不使用的有 \(5\) 人，样本中仅使用 \(A\) 和仅使用 \(B\) 的学生的支付金额分布情况如下：

\begin{center}
\begin{tabular}{c|ccc}
支付方式 & \((0,1000]\) & \((1000,2000]\) & 大于 \(2000\) \\
\hline
仅使用 \(A\) & \(18\) 人 & \(9\) 人 & \(3\) 人 \\
仅使用 \(B\) & \(10\) 人 & \(14\) 人 & \(1\) 人
\end{tabular}
\end{center}

\begin{enumerate}
\item 从全校学生中随机抽取 \(1\) 人，估计该学生上个月 \(A\)，\(B\) 两个支付方式都使用的概率；

\item 从样本仅使用 \(A\) 和仅使用 \(B\) 的学生中各随机抽取 \(1\) 人，以 \(X\) 表示这 \(2\) 人中上个月支付金额大于 \(1000\) 元的人数，求 \(X\) 的分布列和数学期望；

\item 已知上个月样本学生的支付方式在本月没有变化，现从样本仅使用 \(A\) 的学生中，随机抽查 \(3\) 人，发现他们本月的支付金额大于 \(2000\) 元. 根据抽查结果，能否认为样本仅使用 \(A\) 的学生中本月支付金额大于 \(2000\) 元的人数有变化？说明理由.
\end{enumerate}
\end{problem}


\begin{problem}
已知抛物线 \(C:x^2=-2py\) 经过点 \((2,-1)\).

\begin{enumerate}
\item 求抛物线 \(C\) 的方程及其准线方程；

\item 设 \(O\) 为原点，过抛物线 \(C\) 的焦点作斜率不为 \(0\) 的直线 \(l\) 交抛物线 \(C\) 于两点 \(M\)，\(N\)，直线 \(y=-1\) 分别交直线 \(OM\)，\(ON\) 于点 \(A\) 和点 \(B\)，求证：以 \(AB\) 为直径的圆经过 \(y\) 轴上的两个定点.
\end{enumerate}
\end{problem}


\begin{problem}
已知函数 \(f(x)=\dfrac14x^3-x^2+x\).

\begin{enumerate}
\item 求曲线 \(y=f(x)\) 的斜率为 \(1\) 的切线方程；

\item 当 \(x\in[-2,4]\) 时，求证：\(x-6\le f(x)\le x\)；

\item 设 \(F(x)=\left|f(x)-(x+a)\right|\)（\(a\in\R\)），记 \(F(x)\) 在区间 \([-2,4]\) 上的最大值为 \(M(a)\)，当 \(M(a)\) 最小时，求 \(a\) 的值.
\end{enumerate}
\end{problem}


\begin{problem}
已知数列 \(\{a_n\}\)，从中选取第 \(i_1\) 项，第 \(i_2\) 项，…，第 \(i_m\) 项（\(i_1<i_2<\cdots<i_m\)），若 \(a_{i_1}<a_{i_2}<\cdots<a_{i_m}\)，则称新数列 \(a_{i_1}\)，\(a_{i_2}\)，\(\dots\)，\(a_{i_m}\) 为 \(\{a_n\}\) 的长度为 \(m\) 的递增子列. 规定：数列 \(\{a_n\}\) 的任意一项都是 \(\{a_n\}\) 的长度为 \(1\) 的递增子列.

\begin{enumerate}
\item 写出数列 \(1\)，\(8\)，\(3\)，\(7\)，\(5\)，\(6\)，\(9\) 的一个长度为 \(4\) 的递增子列；

\item 已知数列 \(\{a_n\}\) 的长度为 \(p\) 的递增子列的末项的最小值为 \(a_{m_0}\)，长度为 \(q\) 的递增子列的末项的最小值为 \(a_{n_0}\). 若 \(p<q\)，求证：\(a_{m_0}<a_{n_0}\)；

\item 设无穷数列 \(\{a_n\}\) 的各项均为正整数，且任意两项均不相等. 若 \(\{a_n\}\) 的长度为 \(s\) 的递增子列末项的最小值为 \(2s-1\)，且长度为 \(s\) 末项为 \(2s-1\) 的递增子列恰有 \(2^{\,s-1}\) 个（\(s=1\)，\(2\)，\(\dots\)），求数列 \(\{a_n\}\) 的通项公式.
\end{enumerate}
\end{problem}
